\documentclass{article}
\usepackage{amsmath}
\title{A Multi-File Fixture}
\author{Ligature Tests}
\date{}
\begin{document}
\maketitle

\section{Introduction}
This document is split over several files.
The main file holds the introduction and the conclusion,
while the chapters live in a subfolder.
Each sentence sits on its own source line, so that forward and inverse
search can be checked line by line.

Here is a second paragraph with some inline mathematics, $a^2 + b^2 = c^2$,
and a displayed equation:
\begin{equation}
  \int_0^1 x^2 \, dx = \frac{1}{3}
\end{equation}
followed by more text after the equation, which continues the paragraph.

\section{A Chapter in a Subfolder}\label{sec:intro}
This text comes from \texttt{chapters/intro.tex}.
It is read with \verb|\input| from the main file.
Inverse search on these lines must name the subfolder file, not the main file.

A new paragraph in the same subfolder file.
It has \emph{emphasis}, \textbf{bold text} and a box: \fbox{boxed words}.
Some more words to make the paragraph span several lines of output,
so that one source line maps to part of a line and one line of output holds several source lines.

\subsection{Mathematics}
An aligned block:
\begin{align}
  f(x) &= x^2 + 2x + 1 \\
  g(x) &= (x + 1)^2
\end{align}
and a sentence after it.


\section{Lists}
\begin{itemize}
  \item The first item of the list.
  \item A second item that is long enough to wrap onto a second line of output when typeset in the default article layout.
  \item A third item.
\end{itemize}
\begin{enumerate}
  \item Numbered one.
  \item Numbered two, with a footnote.\footnote{This is the footnote text.}
\end{enumerate}

\section{Nested Folder}
This file is two folders deep: \texttt{chapters/details/deep.tex}.
It contains a short paragraph.

\begin{quote}
A quotation, indented on both sides, which is long enough to need more than one line in the output.
\end{quote}

\begin{center}
Centred text on its own.
\end{center}

A last paragraph of the nested file,
written over two source lines.


\section{Conclusion}
The conclusion is back in the main file.
It refers to Section~\ref{sec:intro} and to the table on page~\pageref{tab:small}.
\begin{table}[h]
  \centering
  \begin{tabular}{lcr}
    left & centre & right \\
    alpha & beta & gamma \\
  \end{tabular}
  \caption{A small table.}\label{tab:small}
\end{table}
That is all.
\end{document}
